% fig_collective_screen.tex -- native pgfplots (single column). The moment screen RUN on the collective channels.
% Demonstrated (not asserted) breadth: for BOTH the charge S(q,omega) and spin S^zz(q,omega) structure factors of the
% same Hubbard ring, an m0-preserving corruption (a spurious low-omega phantom peak) passes the total-weight check yet
% moves the first moment. The two channels are OVERLAID in one panel (the L=12 spectra themselves are Fig. sqw);
% here the point is that the SAME screen transports to both. Data: run_collective_sumrule_falsifier.py (L=6, U=8, PBC,
% half filling; exact simulation), representative momentum q/pi=2/3. Operator moments match the exact Lehmann moments
% to <= 6.3e-15 (worst case over q and both channels: summary.worst_lehmann_m1_match_residual in
% data/2026-08-26_collective_sumrule_falsifier.json). The dashed 5% line is a fixed RELATIVE guide, not the deployed
% threshold rule tau_1.
% 2026-09-28 (review pass E: V1, V3, V12, R3-15): the spin channel's m0 residual (column r0 of
% collective_screen_spin.dat) is now drawn as well, so "m0 stays blind for both channels" is shown, not asserted:
% charge m0 markers (squares) at f = 0, 10, ..., 50 and spin m0 markers (open triangles) at f = 5, 15, ..., 45 on the
% same zero line; markers are drawn solid (the dash pattern used to break them); legend words "(flags)" dropped (the
% curves cross a 5% guide, not a rejection threshold); width = \columnwidth like the neighbouring single-column
% figures, height lowered from 6.0 cm to 5.2 cm.
\begin{tikzpicture}
\begin{axis}[
  paperaxis, width=\columnwidth, height=5.2cm,
  axis on top=false,   % grid below data and labels
  xlabel={band weight moved to the spurious peak (\%)},
  ylabel={relative residual $\Delta_1/m_1$ (\%)},
  xmin=0, xmax=50, ymin=-5, ymax=50,
  ytick={0,10,20,30,40,50}, clip=false,
  legend style={at={(0.03,0.97)}, anchor=north west, font=\small, draw=none, fill=none, row sep=-1pt},
  legend cell align=left,
]
  \addplot[threshold, forget plot] coordinates {(0,5) (50,5)};
  \node[pCrimson, font=\scriptsize, anchor=south east, inner sep=1.5pt] at (axis cs:49.5,5.8) {$5\%$ relative guide};
  % charge channel: m1 moves
  \addplot[firesline] table[x=f,y=r1]{figs/collective_screen_charge.dat};
  \addlegendentry{charge $m_1$}
  % spin channel: m1 moves (distinct style); markers solid
  \addplot[pIndigo, very thick, densely dashed, mark=triangle*, mark size=2.4pt, mark options={fill=pIndigo, solid}]
      table[x=f,y=r1]{figs/collective_screen_spin.dat};
  \addlegendentry{spin $m_1$}
  % total weight blind for BOTH channels (m0 is preserved): spin m0 (open triangles at odd points, forget plot) and
  % charge m0 (squares at even points), on the same zero line
  \addplot[pBlind, thick, dashed, mark=triangle*, mark size=2.2pt, mark options={fill=white, draw=pBlind, solid},
           mark repeat=2, mark phase=2, forget plot]
      table[x=f,y=r0]{figs/collective_screen_spin.dat};
  \addplot[blindline, mark options={fill=white, draw=pBlind, solid}, mark repeat=2, mark phase=1,
           legend image code/.code={
             \draw[pBlind, thick, dashed] (0cm,0cm) -- (0.6cm,0cm);
             \draw[pBlind, solid, fill=white] plot[mark=square*, mark size=2pt] coordinates {(0.15cm,0cm)};
             \draw[pBlind, solid, fill=white] plot[mark=triangle*, mark size=2.2pt] coordinates {(0.45cm,0cm)};}]
      table[x=f,y=r0]{figs/collective_screen_charge.dat};
  \addlegendentry{$m_0$, charge \& spin (blind)}
\end{axis}
\end{tikzpicture}
